\section{Finite descent and exponent bounds}\label{sec:towers}
The exponent argument uses two consequences of a finite descent bound.
First, such a bound forces a horizontal vanishing line on a specified
page of the descent spectral sequence. Second, an $E_2$ vanishing line,
together with some finite descent bound, forces a specified fiber map
to have zero image after every chosen exact test. We prove these
statements directly from finite tensor cubes, keeping tested images
distinct from null maps.

\subsection{Finite tensor cubes}\label{par:tower-setup}
Let $\mathcal C$ be a stable symmetric monoidal $\infty$-category whose tensor is
exact in each variable. Let $A$ be an associative unital algebra, with unit
$\one$, and put $J=\fib(\one\to A)$ and $\epsilon:J\to\one$.
For $V\in\mathcal C$ form the augmented Amitsur object
\[
 V\longrightarrow C_A^\bullet(V),\qquad
 C_A^q(V)=A^{\otimes(q+1)}\otimes V.
\]
Cofaces insert the unit and codegeneracies multiply adjacent $A$-factors.
The partial totalization $\Tot_s$ is the limit over the full simplex
category on $[0],\ldots,[s]$, including codegeneracies. The cubical
comparison is standard
descent machinery \cite[Proposition~2.14, pp.~1007--1008]{mnn2017}; its proof fixes the
finite-stage indexing and functoriality used here.

\begin{lemma}[Full finite tensor cube]\label{lem:cube}
For $s\geq0$ there is a natural equivalence of arrows over $\mathrm{id}_V$
in $\mathcal C_{/V}$:
\[
 \begin{aligned}
 &\left(\fib\bigl(V\to\Tot_s C_A^\bullet(V)\bigr)\longrightarrow V\right)\\
 &\qquad\simeq\left(J^{\otimes(s+1)}\otimes V
 \xrightarrow{\epsilon^{\otimes(s+1)}\otimes\mathrm{id}_V}V\right).
 \end{aligned}
\]
The transition on the fibers from stage $s+1$ to stage $s$ applies
$\epsilon$ to the final $J$-factor. Any exact functor
$\Phi:\mathcal C\to\mathrm{Sp}$ preserves this finite fiber sequence.
\end{lemma}
\begin{proof}
Let $\mathcal P_{\ne\varnothing}([s])$ be the poset of nonempty
subsets of $\{0,\ldots,s\}$. The functor
\[
 \theta_s:\mathcal P_{\ne\varnothing}([s])\longrightarrow\Delta_{\leq s},
 \qquad S\longmapsto[|S|-1]
\]
uses the order-preserving injection induced by each inclusion of subsets.
At $[l]$, with $l\leq s$, the comma category $\theta_s/[l]$ consists of
nonempty subsets equipped with a weakly increasing map to $[l]$. It is
the poset of nonempty faces of the simplicial complex $P(s,l)$ whose
simplices are lists
\[
 (i_0,j_0),\ldots,(i_b,j_b),\quad
 0\leq i_0<\cdots<i_b\leq s,\quad
 0\leq j_0\leq\cdots\leq j_b\leq l.
\]
This complex is contractible. The base case $P(0,0)$ is a point.
Every simplex either omits the first-coordinate value $s$ or contains
exactly one vertex $(s,j)$. Begin inductively with the cone of vertex
$(s,l)$ on $P(s-1,l)$.
This cone is contractible even when $l=s$, when its base is not covered
by induction. For each $j<l$, attach the cone of vertex $(s,j)$ along
its base $P(s-1,j)$. This is its exact intersection with the previous
complex and is contractible by induction, since $j\leq s-1$.
The intersections are simplicial subcomplexes, so each attachment
preserves contractibility. The functor is therefore
homotopy initial.

The partial totalization is consequently the limit of the punctured
$(s+1)$-cube obtained by tensoring copies of $\one\to A$ and then
tensoring with $V$. Its omitted vertex is $V$. Taking successive fibers
and using exactness of tensor gives $J^{\otimes(s+1)}\otimes V$,
with its canonical arrow to $V$. Removing the last direction of the
cube gives the stated transition. This is a genuinely finite cubical
limit, so any exact $\Phi$ preserves it. Finiteness of the list of
objects in $\Delta_{\leq s}$ alone would not justify that assertion.
\end{proof}

Fix an exact functor $\Phi:\mathcal C\to\mathrm{Sp}$. Put
$V'=\Phi V$, $U_l=\Phi(J^{\otimes l}\otimes V)$, and
$P_s=\Tot_s(\Phi C_A^\bullet(V))$ for $s\geq0$; set $P_s=0$
for $s<0$. Thus $U_0=V'$ and
\begin{equation}\label{eq:finite-fiber}
 U_{s+1}\longrightarrow V'\longrightarrow P_s
\end{equation}
is a compatible family of fiber sequences. Define its augmentation
filtration by
\begin{equation}\label{eq:augmentation-filtration}
 F^s\pi_dV'=\ker(\pi_dV'\to\pi_dP_{s-1})
           =\im(\pi_dU_s\to\pi_dV'),\qquad s\geq0.
\end{equation}
These groups are defined from finite stages. The argument below will
use them without identifying $V'$ with an infinite totalization.
Our spectral sequence convention is
\[
 E_2^{s,q}=H^s\pi_q(\Phi C_A^\bullet(V)),\qquad
 d_r:(s,q)\longmapsto(s+r,q+r-1).
\]
The full tower has normalized cosimplicial cochains on $E_1$, giving the
displayed $E_2$ page.

\subsection{Exponents and exact finite-page bounds}\label{sec:exponent-pages}
For $b\geq1$ say that $V$ has $A$-exponent at most $b$ if the actual map
\[
 e_b(V)=\epsilon^{\otimes b}\otimes\mathrm{id}_V:
 J^{\otimes b}\otimes V\longrightarrow V
\]
is null. This is stronger than vanishing on homotopy groups after a test.

\begin{lemma}[Exponent calculus]\label{lem:exponent-calculus}
An $A$-module has exponent at most one. Bounds are preserved by retracts
and by tensoring with arbitrary objects. For a cofiber sequence
$U\to V\to W$, bounds $b$ and $c$ on $U$ and $W$ give bound $b+c$ on $V$.
If $e_m(V)$ is null, $V$ is a retract of an object with $m$ successive
layers $A\otimes J^{\otimes i}\otimes V$, $0\leq i<m$.
\end{lemma}
\begin{proof}
The unit map of a module splits by its action, making its fiber map
null. Naturality gives retract invariance, and tensoring a null map
gives tensor invariance. In the cofiber sequence, the composite
$J^{\otimes c}\otimes V\to V\to W$ is null, so choose a lift
$\ell:J^{\otimes c}\otimes V\to U$. After precomposition by $b$
more $\epsilon$ factors, naturality identifies the resulting map into $U$
with $e_b(U)(1\otimes\ell)$, which is null. Its composite to $V$ is
$e_{b+c}(V)$. Finally the consecutive maps
$J^{\otimes(i+1)}\otimes V\to J^{\otimes i}\otimes V$ have the
stated cofibers. Their composites to $V$ filter $\cofib(e_m(V))$ by
these $m$ layers. Nullity of $e_m(V)$ makes $V$ a retract of that cofiber.
\end{proof}

These bounds describe the thick tensor ideal generated by $A$;
compare \cite[Definition~3.18]{mathew2016}. We now fix
the page convention for the quantitative use of nilpotence
\cite{mnn2017,mnn2019}.

\begin{proposition}[Pages and tested images]\label{prop:pages}
For any exact $\Phi:\mathcal C\to\mathrm{Sp}$ and integer $b\geq1$:
\begin{enumerate}
\item If $V$ has $A$-exponent at most $b$, then $E_{b+1}^{s,q}=0$ for all
      $s\geq b$ and $q$.
\item If $V$ has some finite $A$-exponent and $E_2^{s,q}=0$ for all
      $s\geq b$ and $q$, then $\pi_d\Phi(e_b(V))=0$ for every $d$.
\end{enumerate}
The second assertion is zero tested image, not nullity of a map.
\end{proposition}
\begin{proof}
In coordinates $(s,d)$ with $d=q-s$, the derived couple has
\[
 D_r^{s,d}=\im(\pi_dP_{s+r-1}\to\pi_dP_s)
\]
and maps
\[
 i:D_r^{s,d}\to D_r^{s-1,d},\qquad
 j:D_r^{s,d}\to E_r^{s+r,d-1},\qquad
 k:E_r^{s,d}\to D_r^{s,d}.
\]
Here the second $E$-index denotes total degree; no convergence claim is used.
Let $c$ be an actual exponent bound. Every length-$c$ transition in the
fiber tower $U_l$ is null, being $\Phi$ applied to a tensor multiple
of $e_c(V)$. For $s\geq0$ the boundary square from
\eqref{eq:finite-fiber} is
\[
 \begin{array}{ccc}
 \pi_dP_{s+c}&\longrightarrow&\pi_{d-1}U_{s+c+1}\\
 \downarrow&&\downarrow\,0\\
 \pi_dP_s&\longrightarrow&\pi_{d-1}U_{s+1}.
 \end{array}
\]
The image from $P_{s+c}$ has zero boundary and lifts to $\pi_dV'$;
conversely an augmentation class comes from $P_{s+c}$. Hence
\begin{equation}\label{eq:stable-derived-image}
 D_{c+1}^{s,d}=\im(\pi_dV'\to\pi_dP_s).
\end{equation}
Both sides are zero for negative $s$. These image groups have surjective
transition $i$. The augmentation into $P_s$ is injective on homotopy for
$s\geq c-1$, since its kernel is the image of a map factoring through
$e_c(V)$. Thus $i$ is an isomorphism for $s\geq c$. Exactness of the
derived couple gives $j=0$ and, for $s\geq0$, identifies
\[
 E_{c+1}^{s,d}\cong\ker(D_{c+1}^{s,d}\to D_{c+1}^{s-1,d})
                   \cong F^s\pi_dV'/F^{s+1}\pi_dV'.
\]
In particular this is zero for $s\geq c$. Taking $c=b$ proves (i).
For (ii) retain any $c$. The $E_2$ line makes these same quotients zero for
$s\geq b$. Since $F^c=0$, descending from $c$, or monotonicity when
$c<b$, gives $F^b=0$. Equation~\eqref{eq:augmentation-filtration}
is exactly the asserted zero image. This argument uses only finite
cubes and never assumes that $\Phi$ preserves an infinite limit.
\end{proof}
